\documentclass[11pt,a4paper]{article}
\usepackage[utf8]{inputenc}
\usepackage[T1]{fontenc}
\usepackage{amsmath}
\usepackage{amsfonts}
\usepackage{amssymb}
\usepackage{booktabs}
\usepackage{geometry}
\usepackage[numbers]{natbib}
\geometry{margin=1in}

\title{Harmonic Foundations in Species Counterpoint: Roots and Inversions}
\author{}
\date{}

\begin{document}

\maketitle

\section{Introduction}
The application of roots and inversions to the five species of counterpoint represents a synthesis of horizontal melodic independence and vertical harmonic logic. By interpreting the \textit{cantus firmus} not merely as a sequence of intervals but as a part of an underlying fundamental bass, composers achieve a more sophisticated control over musical texture.

\section{Species I: Note-against-Note}
In First Species, where voices move in a 1:1 ratio, the use of inversions is the primary method for avoiding the "jumping" sensation of root-position movement.
\begin{itemize}
    \item \textbf{Consonance and Stability:} While root-position triads provide maximum stability, the first inversion ($6/3$) is essential for maintaining melodic fluidity in the bass \citep{fux1725}.
    \item \textbf{Parallelism:} Using inversions allows for the avoidance of parallel fifths and octaves while maintaining a rich harmonic verticality.
\end{itemize}

\section{Species II and III: Diminution and Arpeggiation}
Species II (2:1) and III (4:1) introduce rhythmic subdivisions. Here, inversions function as tools for arpeggiation.
\begin{itemize}
    \item \textbf{Harmonic Prolongation:} A single root can be sustained across a measure by moving from the root to its inversions (e.g., C to C/E).
    \item \textbf{Passing Tones:} Inversions provide consonant anchor points for dissonant passing tones, ensuring that melodic motion remains anchored to a clear fundamental bass \citep{schoenberg1911}.
\end{itemize}

\section{Species IV: Syncopation and Inverted Resolutions}
Fourth Species focuses on suspensions. The tension created by a suspension is often resolved into an inverted chord to maintain forward momentum.
\begin{itemize}
    \item \textbf{Deferred Finality:} Resolving a suspension to a first inversion triad rather than a root position prevents a premature sense of cadence, keeping the phrase open for further development.
    \item \textbf{Functional Tension:} The use of the second inversion ($6/4$) is strictly controlled here, often appearing as a cadential or passing sonority rather than a stable point of rest \citep{rameau1722}.
\end{itemize}

\section{Species V: Florid Counterpoint and Synthesis}
Species V represents the culmination of counterpoint, combining all previous techniques. The "Rule of the Shortest Way" dictates that voices should transition between various inversions to ensure the smoothest possible voice leading (\textit{Fließender Gesang}). As noted by Schoenberg, the structural function of these inversions is to support the overarching harmonic progression while allowing individual voices to remain melodically interesting \citep{schoenberg1954}.

\begin{thebibliography}{9}
\bibitem{fux1725}
Fux, J. J. (1725). \textit{Gradus ad Parnassum}. (A. Mann, Trans., 1943). W. W. Norton \& Company.

\bibitem{rameau1722}
Rameau, J. P. (1722). \textit{Traité de l'harmonie réduite à ses principes naturels}. Ballard.

\bibitem{schoenberg1911}
Schoenberg, A. (1911). \textit{Theory of Harmony}. (R. E. Carter, Trans., 1978). University of California Press.

\bibitem{schoenberg1954}
Schoenberg, A. (1954). \textit{Structural Functions of Harmony}. Williams and Norgate.
\end{thebibliography}

\end{document}
